\PassOptionsToPackage{table,xcdraw}{xcolor}
\documentclass[11pt, a4paper]{article}

\usepackage[T1]{fontenc}
\usepackage[utf8]{inputenc}
\usepackage{tgpagella}
\usepackage[spanish, es-noshorthands]{babel}
\usepackage{amsmath, amssymb, bm}
\usepackage[most]{tcolorbox}
\usepackage{tabularx}
\usepackage[american]{circuitikz}
\usepackage[margin=2.5cm]{geometry}
\usepackage{fancyhdr}

\pagestyle{fancy}
\fancyhf{}
\setlength{\headheight}{16pt}
\lhead{\textbf{UCB Sede Tarija} | Ing. Mecatrónica}
\rhead{IMT-221 | Guía de Derivación S07-1}
\renewcommand{\headrulewidth}{0.4pt}
\definecolor{ucbblue}{RGB}{0, 51, 102}

\begin{document}

\begin{tcolorbox}[colback=ucbblue!5, colframe=ucbblue, title=\textbf{Guía Asincrónica: Derivación Matemática Extendida de CE, CC y CB}]
    \textbf{Propósito del artefacto:} Esta guía es un complemento de profundización matemática[cite: 14]. Su propósito es permitir reconstruir, paso a paso, las relaciones principales de las tres configuraciones BJT[cite: 14]. La meta no es memorizar tres formularios distintos, sino reconocer que las tres configuraciones nacen del \textbf{mismo modelo incremental}[cite: 14].
\end{tcolorbox}

\section{Convenciones e Hipótesis de Partida}
A temperatura ambiente ($T \approx 300\,\text{K}$), se asume $V_T \approx 25.85\,\text{mV}$[cite: 14]. El modelo incremental se evalúa en el punto de operación $Q$:
\begin{equation*}
    g_m = \frac{I_{CQ}}{V_T}, \qquad r_\pi = \frac{\beta}{g_m}, \qquad r_e \approx \frac{1}{g_m} \text{ (asumiendo $\alpha \approx 1$)} \quad \text{[cite: 14]}
\end{equation*}

\textbf{Hipótesis base para todas las derivaciones[cite: 14]:}
\begin{itemize}
    \item Transistor polarizado en región activa directa (forward-active)[cite: 14].
    \item Análisis de pequeña señal en banda media (capacitores son cortocircuitos)[cite: 14].
    \item Las fuentes DC ideales se reemplazan por \textbf{Tierra AC}[cite: 14].
    \item Se desprecia el Efecto Early ($r_o \to \infty$) para las derivaciones principales[cite: 14].
\end{itemize}

\vspace{0.5cm}
\hrule
\vspace{0.5cm}

\section{Emisor Común (CE)}

\subsection{Topología y Equivalente AC}
\begin{center}
\begin{circuitikz}[scale=1, transform shape, thick]
    \draw (0,0) node[npn](Q){};
    \draw (Q.E) node[ground]{};
    \draw (Q.C) to[R, l_={$R_C$}] (0,2.5) node[ground, rotate=180]{};
    \draw (Q.B) to[short, -o] (-1.5,0) node[left]{$v_i$};
    \draw (Q.C) to[short, *-o] (1.5,0.77) node[right]{$v_o$};
    \node[blue] at (Q.E) [right=0.2cm] {Terminal Común en AC};
\end{circuitikz}
\end{center}
\textbf{Identificación:} Entrada por la base, salida por el colector, emisor común a ambos[cite: 14].

\subsection{Derivación Paso a Paso de la Ganancia ($A_v$)}
Utilizamos el modelo Híbrido-$\pi$[cite: 14].
\begin{enumerate}
    \item Como el emisor está directamente conectado a tierra AC, la tensión base-emisor incremental ($v_\pi$) es exactamente la tensión de entrada[cite: 14]:
    \begin{equation*}
        v_\pi = v_i \quad \text{[cite: 14]}
    \end{equation*}
    \item La fuente de corriente dependiente del modelo extrae corriente desde el colector hacia el emisor[cite: 14]:
    \begin{equation*}
        i_c = g_m v_\pi \implies i_c = g_m v_i \quad \text{[cite: 14]}
    \end{equation*}
    \item En el nodo del colector, la corriente $i_c$ atraviesa la resistencia total de colector $R_C' = R_C \parallel R_L$ (si existe carga) desde tierra hacia el transistor[cite: 14]. Por la ley de Ohm, esto produce una caída de tensión negativa[cite: 14]:
    \begin{equation*}
        v_o = -i_c R_C' \quad \text{[cite: 14]}
    \end{equation*}
    \item Sustituyendo $i_c$:
    \begin{equation*}
        v_o = -(g_m v_i) R_C' \quad \text{[cite: 14]}
    \end{equation*}
    \item Dividiendo entre $v_i$ obtenemos la ganancia:
    \begin{equation*}
        \boxed{A_{v,CE} = \frac{v_o}{v_i} \approx -g_m (R_C \parallel R_L)} \quad \text{[cite: 14]}
    \end{equation*}
\end{enumerate}
\textbf{Interpretación del signo:} Si $v_i \uparrow \implies i_c \uparrow \implies$ la caída en $R_C$ aumenta $\implies$ el voltaje en el colector ($v_o$) disminuye[cite: 14]. Por tanto, \textbf{CE es inversor}[cite: 14].

\subsection{Derivación de Impedancias ($R_{in}$ y $R_{out}$)}
\begin{itemize}
    \item \textbf{Resistencia de Entrada:} Por definición $i_b = v_\pi / r_\pi$[cite: 14]. Como $v_i = v_\pi$, la resistencia vista mirando hacia la base del BJT es $R_{in,BJT} = v_i / i_b = \mathbf{r_\pi}$[cite: 14]. Si hay resistencia de polarización $R_B$ a tierra AC, $R_{in,CE} = \mathbf{R_B \parallel r_\pi}$[cite: 14].
    \item \textbf{Resistencia de Salida:} Para hallar $R_{out}$, anulamos la fuente de entrada ($v_i = 0$)[cite: 14]. Esto causa que $v_\pi = 0$, apagando la fuente dependiente ($g_m v_\pi = 0$)[cite: 14]. La única trayectoria que queda hacia tierra desde la salida es $R_C$ (asumiendo $r_o \to \infty$)[cite: 14]. Por tanto, $R_{out,CE} \approx \mathbf{R_C}$[cite: 14].
\end{itemize}

\vspace{0.5cm}
\hrule
\vspace{0.5cm}

\section{Colector Común (CC) / Seguidor de Emisor}

\subsection{Topología y Equivalente AC}
\begin{center}
\begin{circuitikz}[scale=1, transform shape, thick]
    \draw (0,0) node[npn](Q){};
    \draw (Q.C) -- (0,1.5) node[ground, rotate=180]{};
    \draw (Q.B) to[short, -o] (-1.5,0) node[left]{$v_i$};
    \draw (Q.E) to[R, l_={$R_E$}] (0,-2.5) node[ground]{};
    \draw (Q.E) to[short, *-o] (1.5,-0.77) node[right]{$v_o$};
    \node[blue] at (Q.C) [right=0.2cm] {Terminal Común ($\Delta V_{CC} = 0$)};
\end{circuitikz}
\end{center}
\textbf{Identificación:} Entrada en base, salida en emisor[cite: 14]. ¿Por qué el colector es común? Para una fuente DC ideal $\Delta V_{CC} = 0$, por lo que $V_{CC}$ se convierte en AC ground[cite: 14].

\subsection{Derivación de la Ganancia ($A_v$)}
Sea la carga equivalente del emisor $R_E' = R_E \parallel R_L$[cite: 14].

\textbf{Vía Modelo T (Física del emisor)[cite: 14]:}
\begin{enumerate}
    \item La caída base-emisor es: $v_{be} = i_e r_e$[cite: 14].
    \item La tensión de salida en el emisor es: $v_o = i_e R_E'$[cite: 14].
    \item Por Ley de Voltajes de Kirchhoff, la entrada total es: $v_i = v_{be} + v_o$[cite: 14].
    \item Sustituyendo: $v_i = i_e r_e + i_e R_E' = i_e (r_e + R_E')$[cite: 14].
    \item La ganancia es la razón $v_o / v_i$[cite: 14]:
    \begin{equation*}
        \boxed{A_{v,CC} = \frac{i_e R_E'}{i_e (r_e + R_E')} = \frac{R_E'}{r_e + R_E'}} \quad \text{[cite: 14]}
    \end{equation*}
    Como todas las resistencias son positivas, $0 < A_v < 1$[cite: 14]. Si $R_E' \gg r_e$, entonces $A_{v,CC} \approx 1$[cite: 14].
\end{enumerate}

\textbf{Vía Modelo Híbrido-$\pi$ (Comprobación)[cite: 14]:}
\begin{enumerate}
    \item Sabemos que $i_b = v_\pi / r_\pi$ e $i_e = (\beta + 1)i_b$[cite: 14].
    \item Salida: $v_o = (\beta + 1)i_b R_E'$[cite: 14].
    \item Entrada: $v_i = v_\pi + v_o = i_b r_\pi + (\beta + 1)i_b R_E' = i_b [r_\pi + (\beta + 1)R_E']$[cite: 14].
    \item Dividiendo $v_o / v_i$: $A_v = \frac{(\beta+1)R_E'}{r_\pi + (\beta+1)R_E'}$[cite: 14].
    \item Dividiendo numerador y denominador entre $(\beta+1)$ y recordando que $\frac{r_\pi}{\beta+1} \approx r_e$, recuperamos[cite: 14]:
    \begin{equation*}
        A_v = \frac{R_E'}{r_\pi/(\beta+1) + R_E'} \approx \frac{R_E'}{r_e + R_E'} \quad \text{[cite: 14]}
    \end{equation*}
\end{enumerate}

\subsection{Derivación de $R_{in}$ y $R_{out}$}
\begin{itemize}
    \item \textbf{Resistencia de Entrada ($R_{in}$):} Usando la ecuación de entrada del modelo Híbrido-$\pi$: $v_i = i_b [r_\pi + (\beta+1)R_E']$[cite: 14]. Por definición $R_{in,B} = v_i / i_b$[cite: 14]:
    \begin{equation*}
        \boxed{R_{in,B} = r_\pi + (\beta+1)R_E' \approx (\beta+1)(r_e + R_E')} \quad \text{[cite: 14]}
    \end{equation*}
    \item \textbf{Resistencia de Salida ($R_{out}$):} Anulamos la fuente de entrada. La resistencia de la fuente $R_{sig}'$ vista desde la base refleja su impedancia hacia el emisor dividida por $(\beta+1)$[cite: 14].
    \begin{equation*}
        \boxed{R_{out,CC} \approx R_E \parallel \left( r_e + \frac{R_{sig}'}{\beta+1} \right)} \quad \text{[cite: 14]}
    \end{equation*}
    Si la fuente es ideal ($R_{sig}' \to 0$) y $R_E \gg r_e$, entonces $R_{out,CC} \approx \mathbf{r_e}$ (Muy baja)[cite: 14].
\end{itemize}

\vspace{0.5cm}
\hrule
\vspace{0.5cm}

\section{Base Común (CB)}

\subsection{Topología y Equivalente AC}
\begin{center}
\begin{circuitikz}[scale=1, transform shape, thick]
    \draw (0,0) node[npn](Q){};
    \draw (Q.B) node[ground]{};
    \draw (Q.E) to[short, -o] (-1.5,-0.77) node[left]{$v_i$};
    \draw (Q.C) to[R, l_={$R_C$}] (0,2.5) node[ground, rotate=180]{};
    \draw (Q.C) to[short, *-o] (1.5,0.77) node[right]{$v_o$};
    \node[blue] at (Q.B) [right=0.2cm] {Terminal Común AC};
\end{circuitikz}
\end{center}
\textbf{Identificación:} Entrada en el emisor, salida en colector, base atada a Tierra AC[cite: 14].

\subsection{Derivación de la Ganancia ($A_v$) y $R_{in}$}
\begin{enumerate}
    \item \textbf{Resistencia de entrada:} Al entrar por el emisor (usando el modelo T), la tensión de entrada $v_i$ sostiene directamente la corriente $i_e$ a través de la resistencia intrínseca $r_e$: $v_i \approx i_e r_e$[cite: 14]. 
    \begin{equation*}
        \boxed{R_{in,CB} = \frac{v_i}{i_e} \approx r_e \approx \frac{1}{g_m}} \quad \text{[cite: 14]}
    \end{equation*}
    Esto explica por qué la CB tiene una resistencia de entrada naturalmente bajísima[cite: 14].
    
    \item \textbf{Ganancia de tensión:} Como la base está en tierra AC, $v_b = 0$[cite: 14].
    \item $v_{be} = v_b - v_e = 0 - v_i = -v_i$[cite: 14].
    \item La fuente de corriente del colector es $i_c = g_m v_{be}$[cite: 14]. Sustituyendo: $i_c = -g_m v_i$[cite: 14].
    \item La salida en el colector es $v_o = -i_c R_C'$ (con $R_C' = R_C \parallel R_L$)[cite: 14].
    \item Sustituyendo $i_c$: $v_o = -(-g_m v_i) R_C' = g_m R_C' v_i$[cite: 14].
    \begin{equation*}
        \boxed{A_{v,CB} = \frac{v_o}{v_i} \approx +g_m (R_C \parallel R_L)} \quad \text{[cite: 14]}
    \end{equation*}
    El signo positivo demuestra que \textbf{CB no invierte} la fase de tensión[cite: 14].
\end{enumerate}
\textbf{Nota sobre corriente:} La ganancia de corriente es $A_i = i_c / i_e = \alpha = \frac{\beta}{\beta+1} \approx 1$[cite: 14]. Por tanto, transfiere la corriente casi íntegramente mientras amplifica la tensión gracias a que la salida pasa por un gran resistor $R_C$[cite: 14].

\newpage

\section{Comparación Matemática Unificada}
\renewcommand{\arraystretch}{1.8}
\begin{tabularx}{\textwidth}{|>{\raggedright\arraybackslash}p{3cm}|>{\raggedright\arraybackslash}X|>{\raggedright\arraybackslash}X|>{\raggedright\arraybackslash}X|}
\hline
\rowcolor{ucbblue!20}
\textbf{Cantidad} & \textbf{Emisor Común (CE)} & \textbf{Colector Común (CC)} & \textbf{Base Común (CB)} \\ \hline
\textbf{Entrada} & Base & Base & Emisor \quad \text{[cite: 14]} \\ \hline
\textbf{Salida} & Colector & Emisor & Colector \quad \text{[cite: 14]} \\ \hline
\textbf{Común AC} & Emisor & Colector & Base \quad \text{[cite: 14]} \\ \hline
\textbf{Ganancia $A_v$} & $-g_m R_C'$ & $\frac{R_E'}{r_e + R_E'}$ & $+g_m R_C'$ \quad \text{[cite: 14]} \\ \hline
\textbf{Fase} & Inversora & No inversora & No inversora \quad \text{[cite: 14]} \\ \hline
\textbf{$R_{in}$ (transistor)} & $r_\pi$ & $r_\pi + (\beta+1)R_E'$ & $r_e$ \quad \text{[cite: 14]} \\ \hline
\textbf{$R_{out}$ tendencia} & Moderada/Alta & Baja & Moderada/Alta \quad \text{[cite: 14]} \\ \hline
\textbf{Rol primario} & Ganancia de tensión & Buffer & Interfaz baja-$R_{in}$ \quad \text{[cite: 14]} \\ \hline
\end{tabularx}

\section{Mismo Punto Q, Tres Funciones}
Suponga $I_C = 1\,\text{mA}$, $\beta=100$, $T=300\,\text{K}$. Obtenemos $g_m \approx 38.68\,\text{mS}$, $r_e \approx 25.85\,\Omega$, $r_\pi \approx 2.585\,\text{k}\Omega$[cite: 14].
\begin{itemize}
    \item \textbf{CE (con $R_C'=2\,\text{k}\Omega$):} $A_v \approx -(0.03868)(2000) = \mathbf{-77.4}$. \textit{Gran ganancia invertida}[cite: 14].
    \item \textbf{CC (con $R_E'=1\,\text{k}\Omega$):} $A_v = \frac{1000}{1000 + 25.85} \approx \mathbf{0.975}$. $R_{in,B} \approx 101(1025.85) = \mathbf{103.6\,\text{k}\Omega}$. \textit{Ganancia cercana a 1 y alta $R_{in}$}[cite: 14].
    \item \textbf{CB (con $R_C'=2\,\text{k}\Omega$):} $A_v \approx +(0.03868)(2000) = \mathbf{+77.4}$. $R_{in} \approx \mathbf{25.9\,\Omega}$. \textit{Gran ganancia no invertida y baja $R_{in}$}[cite: 14].
\end{itemize}

\section{Preguntas de Autoevaluación}
\begin{enumerate}
    \item ¿Por qué el Emisor Común (CE) es inversor y la Base Común (CB) no lo es? (Sugerencia: Siga la ecuación $v_o = -i_c R_C'$ en ambos modelos)[cite: 14].
    \item ¿Por qué el Colector Común refleja la resistencia del emisor hacia la base multiplicada por aproximadamente $(\beta+1)$?[cite: 14]
    \item ¿Cómo puede la Base Común tener una ganancia de corriente $A_i \approx 1$ y al mismo tiempo una ganancia de tensión $A_v \gg 1$? (Sugerencia: Analice la diferencia de resistencias donde la corriente ingresa vs. donde sale)[cite: 14].
    \item ¿Por qué $r_o \to \infty$ es una hipótesis de simplificación y no una propiedad real del transistor?[cite: 14]
\end{enumerate}

\section{Flujo de Estudio Recomendado}
Para dominar cualquier configuración sin memorizar fórmulas[cite: 14]:
\begin{enumerate}
    \item Redibujar el circuito identificando entrada, salida y terminal común[cite: 14].
    \item Construir el equivalente AC y reemplazar el transistor por el modelo Híbrido-$\pi$ o T[cite: 14].
    \item Derivar $A_v = v_o/v_i$ paso a paso mediante Ley de Ohm y Kirchhoff[cite: 14].
    \item Derivar $R_{in} = v_i/i_i$ analizando el puerto de entrada[cite: 14].
    \item Solo al final, comparar los resultados propios con la tabla unificada[cite: 14].
\end{enumerate}

\end{document}
